\chapter{Curve Representation of Extracted Topology}

% Caption setup for algorithm 
\captionsetup[algorithm]{
	labelfont=bf,
	justification=raggedright,
	singlelinecheck=false
}

%% ???
Curve representation of the extracted topology is performed in three main steps. First, the topology pixels are converted into polylines. Next, redundant points are removed. Finally, the resulting polylines are converted into curves.

\section{Problem Formulation}

\section{Preliminaries}

\section{Endpoint and Junction Pixel Detection}

One of the key steps in creating polylines from topological information is determining whether a pixel is an endpoint or a junction. We use a simple classification method based on counting the number of foreground pixels in the 8-neighborhood of the examined pixel.

If the number of foreground neighbors is one, the pixel is classified as an endpoint, since it is connected to only one neighboring pixel. If the number is two, the pixel is considered a regular polyline pixel connecting two neighboring pixels.

% todo obrazek endpointu, normální, junction, a abomination jako plateau or other.

A junction pixel has three or more foreground neighbors, although in typical cases it has exactly three. If four or more foreground neighbors are present, this may indicate a more complex local structure, such as a plateau or another shape that can complicate polyline extraction.

\section{Polyline Creation Algorithm}

The polyline creation algorithm uses topological information, such as a binary mask skeleton or contour, to construct a continuous polyline that covers all pixels present in the topology. The algorithm relies on the previously detected endpoint and junction pixels and proceeds in three main steps.

\subsection{Polyline from Endpoint to Junction or Endpoint}

The first step is to create polylines starting at an endpoint and ending either at a junction or at another endpoint. During the process, we keep track of visited pixels as well as all detected endpoints and junctions. The algorithm starts by adding the initial endpoint to the polyline and marking it as visited. It then searches for foreground pixels in its 8-neighborhood. 

If a foreground neighbor is found and it has already been visited, it is skipped. Otherwise, the algorithm determines whether the pixel is an endpoint, a junction, or a regular pixel and performs one of the following three actions:

\begin{itemize}
	\item If the pixel is an endpoint, it is marked as visited, added to the polyline, and the resulting polyline is returned.
	\item If the pixel is a junction, it is added to the polyline and the polyline is returned, but the junction pixel itself is not marked as visited.
	\item If the pixel is neither an endpoint nor a junction, it is added to the polyline, marked as visited, and the search continues from that pixel.
\end{itemize}

To improve the accuracy and continuity of the resulting polyline, the endpoint or junction neighbor is prioritized over other foreground pixels, algorithm chooses the closed endpoint or junction neighbor. In our implementation, Manhattan distance was used because of its low computational cost. Without this selection step, the resulting polyline may become discontinuous or may contain short redundant segments.

% todo obrázek bez a s seleckcí, jak vypadá pak polyline

{
\captionof{algorithm}{Polyline from Endpoint to Junction or Endpoint algorithm}
\begin{algorithmic}[1]
	\Procedure{Endpoint\_to\_junction\_or\_endpoint}{$\textit{starting\_endpoint},\textit{visited}$}
	\State $\textit{queue} \gets \text{initialize queue}$
	\State $\textit{polyline} \gets \text{initialize empty polyline}$
	\State
	\State $\textit{queue}\text{.add(}\textit{starting\_endpoint}\text{)}$
	\While{$\textit{queue}$ is not empty}
		\State $\textit{point} \gets \textit{queue}\text{.pop()}$ 
		\State
		\State $\textit{visited}\text{[}\textit{ point }\text{]} = \text{true}$
		\State $\textit{polyline}\text{.add(}\textit{point}\text{)}$
		\State
		\State $\textit{best\_endpoint\_or\_junction}\gets\textbf{ select from neighbors of }\textit{point}$
		\If{$\textit{best\_endpoint\_or\_junction}\text{ exists}$}
			\State $\textit{polyline}\text{.add(}\textit{best\_endpoint\_or\_junction}\text{)}$
			\State
			\Return $\textit{polyline}$
		\EndIf
		\State
		\For {$\textit{neighbor}\textbf{ from }\textit{neighbors}\textbf{ of }\textit{point}$}
			\If{$\textit{neighbor}\text{ is foreground} \textbf{ and } \text{not }\textit{visited}\text{[}\textit{ neighbor }\text{]}$}
				\State $\textit{queue}\text{.add(}\textit{neighbor}\text{)}$
				\State $\textbf{break}$
			\EndIf
		\EndFor
	\EndWhile
	
	\State
	\State
	\Return $\textit{polyline}$
	\EndProcedure
\end{algorithmic}
}

% todo obrázek endpoint - endpoint/junction a vyrobená polyline. (2x)

\subsection{Polyline from Junction to Junction}

The second step is similar to the first one, but instead of constructing a single polyline from an endpoint, it constructs multiple polylines from a given junction pixel. This is expected, since an endpoint has only one adjacent foreground pixel, whereas a junction has at least two. The algorithm follows a procedure similar to that of the previous step. It starts by considering the foreground pixels in the 8-neighborhood of the starting junction as initial neighbor pixels.

If an initial neighbor pixel has already been visited, it is skipped. The number of constructed polylines is equal to the number of unvisited foreground neighbors of the starting junction. Then, the starting junction and the selected neighbor are added to the polyline. If the starting neighbor is itself another junction, it is added to the polyline, and the algorithm continues with the next starting neighbor.

Otherwise, the starting neighbor is marked as visited, and the algorithm continues by searching for foreground pixels in its 8-neighborhood. If a foreground neighbor is found and it has already been visited, it is skipped. Otherwise, the algorithm determines whether the pixel is a junction or a regular pixel and performs one of the following two actions:

\begin{itemize}
	\item If the pixel is a junction, it is added to the polyline, and the algorithm continues with the next starting neighbor by repeating the steps described above.
	\item If the pixel is not a junction, it is added to the polyline, marked as visited, and the search continues from that pixel.
\end{itemize}

To improve the accuracy and continuity of the resulting polylines, junction pixels are prioritized over regular foreground pixels whenever possible.This is achieved in the same way as in the previous step, by selecting the closest junction candidate. In our implementation, we used Manhattan distance because of its low computational cost.

{
\captionof{algorithm}{Polyline from Junction to Junction algorithm}
\begin{algorithmic}[1]
	\Procedure{Junction\_to\_junction}{$\textit{starting\_junction},\textit{visited}$}
	\State $\textit{polyline\_list} \gets \text{initialize empty polyline list}$
	\State
	\State $\textit{visited}\text{[ }\textit{starting\_junction}\text{ ]}$ = true
	\State
	\For {$\textit{starting\_neighbor}\textbf{ from }\textit{neighbors}\textbf{ of }\textit{starting\_junction}$}
		\If{$\textit{visited}\text{[ }\textit{starting\_neighbor}\text{ ]}$ } 
			\State $\textbf{continue}$
		\EndIf
		\State
		\State $\textit{queue} \gets \text{initialize queue}$
		\State $\textit{polyline} \gets \text{initialize empty polyline}$
		\State $\textit{polyline}\text{.add(}\textit{starting\_junction}\text{)}$
		\State
		\State $\textit{queue}\text{.add(}\textit{starting\_neighbor}\text{)}$
		\While{$\textit{queue}$ is not empty}
			\State $\textit{point} \gets \textit{queue}\text{.pop()}$ 
			\State
			\If{$\textit{point}\textbf{ is junction}$}
				\State $\textit{polyline}\text{.add(}\textit{point}\text{)}$
				\State $\textbf{break}$
			\EndIf
			\State
			\State $\textit{visited}\text{[}\textit{ point }\text{]} = \text{true}$
			\State $\textit{polyline}\text{.add(}\textit{point}\text{)}$
			\State
			\State $\textit{best\_endpoint\_or\_junction}\gets\textbf{ select from neighbors of }\textit{point}$
			\If{$\textit{best\_endpoint\_or\_junction}\text{ exists}$}
			\State $\textit{polyline}\text{.add(}\textit{best\_endpoint\_or\_junction}\text{)}$
			\State $\textbf{break}$
			\EndIf
			\State
			\For {$\textit{neighbor}\textbf{ from }\textit{neighbors}\textbf{ of }\textit{point}$}
				\If{$\textit{neighbor}\text{ is foreground} \textbf{ and } \text{not }\textit{visited}\text{[}\textit{ neighbor }\text{]}$}
					\State $\textit{queue}\text{.add(}\textit{neighbor}\text{)}$
					\State $\textbf{break}$
				\EndIf
			\EndFor
		\EndWhile
		\State
		\State $\textit{polyline\_list}\text{.add(}\textit{polyline}\text{)}$
	\EndFor
	\State
	\State
	\Return $\textit{polyline\_list}$
	\EndProcedure
\end{algorithmic}
}

% todo obrázek s resultem junction to junction.

\subsection{Polyline from Closed Contour}

After the first two steps, all remaining unvisited foreground pixels are expected to form closed polylines without endpoints or junctions. For this reason, the algorithm uses the map of previously visited foreground pixels to correctly construct the closed polylines. Closed polylines typically appear when processing contours rather than skeletons.

The algorithm starts by selecting the first unvisited foreground pixel as the starting point. The starting point is then added to the polyline and marked as visited. It then searches for foreground pixels in its 8-neighborhood. 

If an unvisited foreground neighbor is found, it is added to the polyline, marked as visited, and the search continues from that pixel. If no unvisited foreground neighbor exists, the construction of the closed polyline is complete. Finally, the initial starting point is added to the polyline again, since in our implementation a closed polyline is represented by having the same first and last point.

{
	\captionof{algorithm}{Polyline from Closed Contour algorithm}
	\begin{algorithmic}[1]
		\Procedure{Closed\_contour\_to\_polyline}{$\textit{starting\_point},\textit{visited}$}
		\State $\textit{queue} \gets \text{initialize queue}$
		\State $\textit{polyline} \gets \text{initialize empty polyline}$
		\State
		\State $\textit{queue}\text{.add(}\textit{starting\_point}\text{)}$
		\While{$\textit{queue}$ is not empty}
		\State $\textit{point} \gets \textit{queue}\text{.pop()}$ 
		\State
		\State $\textit{visited}\text{[}\textit{ point }\text{]} = \text{true}$
		\State $\textit{polyline}\text{.add(}\textit{point}\text{)}$
		\State
		\For {$\textit{neighbor}\textbf{ from }\textit{neighbors}\textbf{ of }\textit{point}$}
		\If{$\textit{neighbor}\text{ is foreground} \textbf{ and } \text{not }\textit{visited}\text{[}\textit{ neighbor }\text{]}$}
		\State $\textit{queue}\text{.add(}\textit{neighbor}\text{)}$
		\State $\textbf{break}$
		\EndIf
		\EndFor
		\EndWhile
		
		\State
		\State $\textit{polyline}\text{.add(}\textit{starting\_point}\text{)}$
		\State 
		\Return $\textit{polyline}$
		\EndProcedure
	\end{algorithmic}
}

% todo obrázek s udělanou closed polyline

\subsection{Edge Cases}

There are several edge cases that can be handled by the previously described algorithms, but not in an optimal way. In this application, which also relies heavily on user input, we decided not to handle these cases automatically, as doing so may produce unwanted results. Instead, we provide tools that allow the user to resolve such cases manually, for example by stitching two or more polylines together or by closing a polyline.

Examples of such edge cases include the previously mentioned junction plateaus and L-shaped corners. Although the algorithm can process these cases, it may create redundant line segments or unnecessary connections. Another encountered edge case arises when a 2-pixel-wide contour is created from a foreground region that is too thin. This may result in many short line segments and in a polyline that is not closed.

% todo obrázek L, plateau atd a jak to algoritmus řeší -> tedy asi barvevné spojnice středů?

% todo obrázek rozsypané kontůry

\section{Polyline Modification Algorithms}



\subsection{Polyline Stitching Algorithm}



\subsection{Douglas--Peucker Algorithm} % todo check jména, not sure about it :D



\section{Spline Construction from Polylines}

\subsection{Hermite and B\'ezier Curve Representations}

In our implementation, we chose to represent splines as a sequence of cubic B\'ezier curves rather than as Hermite curves, which are used internally by Unreal Engine~5. The main reason for this choice is that B\'ezier curves are more convenient for the algorithms introduced later in this chapter. At the same time, a cubic B\'ezier curve can be converted to a Hermite representation in a straightforward manner.

A cubic B\'ezier curve $\mathbf{B}(t)$ is defined by two end points, $\mathbf{P}_0$ and $\mathbf{P}_3$, and two control points, $\mathbf{P}_1$ and $\mathbf{P}_2$. It is given by
\begin{equation}
	\mathbf{B}(t) = (1 - t)^3\mathbf{P}_0 + 3(1-t)^2t\mathbf{P}_1 + 3(1-t)t^2\mathbf{P}_2 + t^3\mathbf{P}_3,
\end{equation}
for $t \in [0,1]$.

A Hermite curve $\mathbf{H}(t)$ is defined by two end points, $\mathbf{H}_0$ and $\mathbf{H}_1$, and two tangent vectors, $\mathbf{M}_0$ and $\mathbf{M}_1$. It is given by
\begin{equation}
	\mathbf{H}(t) = (2t^3 - 3t^2 + 1)\mathbf{H}_0 + (t^3 - 2t^2 + t)\mathbf{M}_0 + (-2t^3 + 3t^2)\mathbf{H}_1 + (t^3 - t^2)\mathbf{M}_1,
\end{equation}
for $t \in [0,1]$.

A Hermite curve defined by end points $\mathbf{H}_0$, $\mathbf{H}_1$ and tangent vectors $\mathbf{M}_0$, $\mathbf{M}_1$ can be constructed from a cubic B\'ezier curve defined by control points $\mathbf{P}_0$, $\mathbf{P}_1$, $\mathbf{P}_2$, and $\mathbf{P}_3$ using the following relationships:
\begin{align}
	\mathbf{H}_0 &= \mathbf{P}_0, \\
	\mathbf{M}_0 &= 3\left(\mathbf{P}_1 - \mathbf{P}_0\right), \\
	\mathbf{M}_1 &= 3\left(\mathbf{P}_3 - \mathbf{P}_2\right), \\
	\mathbf{H}_1 &= \mathbf{P}_3.
\end{align}

\subsection{Catmull--Rom Spline Construction}

To convert a polyline into a spline represented by B\'ezier curves, we employ a Catmull--Rom spline as an intermediate representation. A Catmull--Rom spline can be viewed as a special case of a cubic Hermite spline in which tangent vectors are determined from neighboring control points.

For the parameterization, we use the centripetal Catmull--Rom spline. Let the knot sequence satisfy
\begin{equation}
	t_{k+1} = t_k + \left\lVert \mathbf{P}_{k+1} - \mathbf{P}_k \right\rVert^\alpha,
\end{equation}
where $\alpha \in [0,1]$. In our implementation, we use $\alpha = 0.5$. This variant was selected because the centripetal Catmull--Rom spline is the only variant that guarantees no cusps or self-intersections within a single curve segment~\cite{bk_catmull_rom_bezier}.

% todo obrázek jak vypadá Catmull--Rom spline mezi čtyřmi body

As follows from the definition, a Catmull--Rom spline segment is constructed from four points, $\mathbf{P}_0$ to $\mathbf{P}_3$, and produces a curve segment between $\mathbf{P}_1$ and $\mathbf{P}_2$. For an open polyline, this means that the spline would not naturally interpolate the first and the last point of the polyline. To resolve this issue, we introduce auxiliary end points defined as
\begin{align}
	\mathbf{P}_{-1} &= 2\mathbf{P}_0 - \mathbf{P}_1, \\
	\mathbf{P}_{n+1} &= 2\mathbf{P}_n - \mathbf{P}_{n-1},
\end{align}
where the polyline consists of points $\mathbf{P}_0, \dots, \mathbf{P}_n$.

% todo obrázek jak vypadá Catmull--Rom spline s doplněnými konci (+ vidět ty nové krajní body)

A drawback of this endpoint extension is that it may produce undesirable shapes when the curve is very sharp near the boundary. However, since the base polylines are constructed from binary mask skeletons or contours, this situation occurs only rarely in practice.

% todo obrázek tohoto edge case

The B\'ezier curve representation of a Catmull--Rom spline segment can be computed directly from the four defining points, $\mathbf{P}_0$ to $\mathbf{P}_3$. If the B\'ezier curve is defined by control points $\mathbf{B}_0$ to $\mathbf{B}_3$, the relationship is given by~\cite{bk_catmull_rom_bezier}
\begin{align}
	\mathbf{B}_0 &= \mathbf{P}_1, \\
	\mathbf{B}_1 &= \frac{d_1^{2\alpha}\mathbf{P}_2 - d_2^{2\alpha}\mathbf{P}_0 + \left(2d_1^{2\alpha} + 3d_1^\alpha d_2^\alpha + d_2^{2\alpha}\right)\mathbf{P}_1}{3d_1^\alpha\left(d_1^\alpha + d_2^\alpha\right)}, \\
	\mathbf{B}_2 &= \frac{d_3^{2\alpha}\mathbf{P}_1 - d_2^{2\alpha}\mathbf{P}_3 + \left(2d_3^{2\alpha} + 3d_3^\alpha d_2^\alpha + d_2^{2\alpha}\right)\mathbf{P}_2}{3d_3^\alpha\left(d_3^\alpha + d_2^\alpha\right)}, \\
	\mathbf{B}_3 &= \mathbf{P}_2,
\end{align}
where $d_1$, $d_2$, and $d_3$ are defined as
\begin{align}
	d_1 &= \left\lVert \mathbf{P}_1 - \mathbf{P}_0 \right\rVert, \\
	d_2 &= \left\lVert \mathbf{P}_2 - \mathbf{P}_1 \right\rVert, \\
	d_3 &= \left\lVert \mathbf{P}_3 - \mathbf{P}_2 \right\rVert,
\end{align}
and $\left\lVert \cdot \right\rVert$ denotes the Euclidean norm.

\subsection{Curve Subdivision Using de Casteljau's Algorithm}